{\rm yr}^{-1}$ \cite{CRflux}, if the observed
flux of ultra-high-energy cosmic-rays is produced by 
sources that are cosmologically distributed. 
The energy-dependent generation rate of cosmic-rays
is therefore given by 
\begin{equation}
E_{CR}^2{d\dot N_{CR}\over dE_{CR}}={\dot\varepsilon_{CR}^
{[10^{19},10^{21}]}\over
\ln(10^{21}/10^{19})}\approx
10^{44}{\rm erg\ Mpc}^{-3}{\rm yr}^{-1}.
\label{ECR}
\end{equation} 
If the high-energy
protons produced by the extra-galactic sources lose a fraction
$\epsilon<1$ of their
energy through photo-meson production of pions before escaping the source, 
the resulting present-day energy density of muon neutrinos is 
$E_\nu^2 dN_\nu/ dE_\nu
\approx0.25\epsilon t_H E_{CR}^2d\dot N_{CR}/dE_{CR}$, 
where $t_H\approx10^{10}{\rm yr}$ is the Hubble time.
For energy independent $\epsilon$ 
the neutrino spectrum follows the proton generation spectrum, since
the fraction of the proton energy carried by a neutrino produced through a
photo-meson interaction, $E_\nu\approx0.05E_p$, 
is independent of the proton energy.
The $0.25$ factor arises because neutral pions, 
which do not produce neutrinos, are produced with roughly equal
probability with charged pions, and because in the decay 
$\pi^+\rightarrow\mu^++\nu_\mu
\rightarrow e^++\nu_e+\overline\nu_\mu+\nu_\mu$
muon neutrinos carry approximately half
the charged pion energy. Defining $I_{\rm max}$ as the muon neutrino intensity 
($\nu_\mu$ and $\bar\nu_\mu$ combined) obtained for $\epsilon=1$, 
\begin{equation}
I_{\rm max}\approx0.25\xi_Zt_H{c\over4\pi}E_{CR}^2{d\dot N_{CR}\over dE_{CR}}
\approx1.5\times10^{-8}\xi_Z{\rm GeV\,cm}^{-2}{\rm s}^{-1}
{\rm sr}^{-1}\,,
\label{Fmax}
\end{equation}
the expected neutrino intensities are
\begin{equation}
E_\nu^2\Phi_{\nu_\mu}\equiv{c\over4\pi} E_\nu^2{dN_{\nu_\mu}\over dE_\nu}
={1\over2}\epsilon I_{\rm max},\quad 
\Phi_{\nu_e}\approx\Phi_{\bar\nu_\mu}\approx\Phi_{\nu_\mu}.
\label{Jnu}
\end{equation}
% ...
of redshift and cosmology; $g(z)\equiv1$ 
for a flat universe with zero cosmological constant. Let us assume that
the neutrino energy generation rate evolves rapidly with redshift, following
the luminosity density evolution of QSOs \cite{QSO}, which may be described
as $f(z)=(1+z)^{\alpha}$ with $\alpha\approx3$ \cite{QSOl} at low redshift, 
$z<1.9$, $f(z)={\rm Const.}$ for $1.9<z<2.7$, and an exponential decay at 
$z>2.7$ \cite{QSOh}. Using this functional form of $f(z)$, which is 
also similar \cite{QSO} to that 
describing the evolution of star formation rate \cite{SFR}, we find that
$\xi_Z\approx3$ (with weak dependence on cosmology). For no evolution,
$f(z)={\rm Const.}$, we have $\xi_Z\approx0.6<1$ 
(with only a weak dependence on cosmology) due to a redshift energy loss of
neutrinos.
